\documentclass[border=10pt]{standalone}
\usepackage{ctex,tikz}
% 引入 positioning 库以便使用 "above of=", "right of=" 等简单定位语法
\usetikzlibrary{positioning}

\begin{document}

\begin{tikzpicture}[
    % 全局样式定义
    node distance = 1.5cm and 1.2cm, % 节点之间的垂直和水平距离
    box/.style = {
        draw, 
        rounded corners, 
        thick, 
        align = center, 
        font = \sffamily\small,
        minimum width = 2.2cm,
        minimum height = 0.8cm
    },
    % mainNode/.style = {
    %     fill = blue!15, 
    %     text = black, 
    %     font = \sffamily\bfseries\large, 
    %     line width = 2pt,
    %     minimum size = 2.5cm,
    %     circle % 中心设为圆形
    % },
    mainNode/.style = {
        box, 
        fill = blue!10, 
        text width = 2.5cm, 
        font = \sffamily\bfseries\normalsize,
        line width = 1.5pt
    },
    subNode/.style = {
        box, 
        fill = white,
        % text width = 1.8cm
        font = \sffamily\small,
    },
    arrow/.style = {
        ->, 
        thick, 
        >=stealth
    }
]

    % --- 1. 核心节点 ---
    \node[mainNode] (center) {深度学习 \\ Deep Learning};

    % --- 2. 上方分支：核心架构 ---
    \node[box, fill=red!10, above=of center] (archTitle) {核心架构};
    \node[subNode, above left=of archTitle] (cnn) {CNN \\ 卷积网络};
    \node[subNode, above=of archTitle] (rnn) {RNN/LSTM \\ 循环网络};
    \node[subNode, above right=of archTitle] (trans) {Transformer \\ 注意力机制};
    
    % --- 3. 右侧分支：应用领域 ---
    \node[box, fill=green!10, right=of center] (appTitle) {应用领域};
    \node[subNode, above right=of appTitle] (cv) {计算机视觉 \\ CV};
    \node[subNode, right=of appTitle] (nlp) {自然语言处理 \\ NLP};
    \node[subNode, below right=of appTitle] (rec) {音频处理 \\ Audio AI};

    % --- 4. 下方分支：优化训练 ---
    \node[box, fill=orange!10, below=of center] (optTitle) {优化与训练};
    \node[subNode, below left=of optTitle] (bp) {反向传播};
    \node[subNode, below=of optTitle] (adam) {优化器 \\ Adam/SGD};
    \node[subNode, below right=of optTitle] (drop) {Dropout \\ 正则化};

    % --- 5. 左侧分支：数学基础 ---
    \node[box, fill=purple!10, left=of center] (mathTitle) {数学基础};
    \node[subNode, above left=of mathTitle] (linear) {线性代数};
    \node[subNode, left=of mathTitle] (calc) {微积分};
    \node[subNode, below left=of mathTitle] (prob) {概率统计};

    % --- 绘制连线 ---
    % 中心到四大分支标题
    \draw[arrow] (center) -- (archTitle);
    \draw[arrow] (center) -- (appTitle);
    \draw[arrow] (center) -- (optTitle);
    \draw[arrow] (center) -- (mathTitle);

    % 分支标题到子节点 (上方)
    \draw[arrow] (archTitle) -- (cnn);
    \draw[arrow] (archTitle) -- (rnn);
    \draw[arrow] (archTitle) -- (trans);

    % 分支标题到子节点 (右侧)
    \draw[arrow] (appTitle) -- (cv);
    \draw[arrow] (appTitle) -- (nlp);
    \draw[arrow] (appTitle) -- (rec);

    % 分支标题到子节点 (下方)
    \draw[arrow] (optTitle) -- (bp);
    \draw[arrow] (optTitle) -- (adam);
    \draw[arrow] (optTitle) -- (drop);

    % 分支标题到子节点 (左侧)
    \draw[arrow] (mathTitle) -- (linear);
    \draw[arrow] (mathTitle) -- (calc);
    \draw[arrow] (mathTitle) -- (prob);

\end{tikzpicture}

\end{document}